\chapter{Mathematics}
\label{ch:maths}

\section{Inline and displayed}\label{inline-and-displayed}

Inline maths sits within a sentence, and displayed maths gets its own line. Load \texttt{amsmath} for anything beyond the basics.

\begin{lstlisting}[language=latex]
The energy is $E = mc^2$ here.

\[ E = mc^2 \]          % displayed, no number

\begin{equation}        % displayed and numbered
  E = mc^2
  \label{eq:energy}
\end{equation}
\end{lstlisting}

In the file below, these forms and some of the building blocks come together. Only the numbered equations get a number at the right.

\textbf{maths.tex}

\lstinputlisting{../../../teaching/latex/examples/maths.tex}

\begin{figure}
\centering
\includegraphics[width=\linewidth]{img/maths.png}
\caption{What that file compiles to.}
\end{figure}

\section{The building blocks}\label{the-building-blocks}

{\def\LTcaptype{none} % do not increment counter
\begin{longtable}[]{@{}ll@{}}
\toprule\noalign{}
You want & You write \\
\midrule\noalign{}
\endhead
\bottomrule\noalign{}
\endlastfoot
Superscript and subscript & \texttt{x\^{}2}, \texttt{x\_i}, \texttt{x\_\{i,j\}} \\
Fraction & \texttt{\textbackslash{}frac\{a\}\{b\}} \\
Square root & \texttt{\textbackslash{}sqrt\{x\}} \\
Greek letters & \texttt{\textbackslash{}alpha}, \texttt{\textbackslash{}mu}, \texttt{\textbackslash{}Delta} \\
Sum and integral & \texttt{\textbackslash{}sum\_\{i=1\}\^{}\{n\}}, \texttt{\textbackslash{}int\_0\^{}1} \\
Operators and relations & \texttt{\textbackslash{}times}, \texttt{\textbackslash{}approx}, \texttt{\textbackslash{}leq}, \texttt{\textbackslash{}rightarrow} \\
Words in maths & \texttt{\textbackslash{}text\{hours\}} \\
Spacing & \texttt{\textbackslash{},} small, \texttt{\textbackslash{}quad} wide \\
Brackets that grow & \texttt{\textbackslash{}left(\ \textbackslash{}frac\{a\}\{b\}\ \textbackslash{}right)} \\
\end{longtable}
}

Braces group things, so \texttt{x\^{}\{10\}} raises the whole 10, while \texttt{x\^{}10} raises only the 1.

\section{Aligned derivations}\label{aligned-derivations}

Long derivations read best one step per line. Use \texttt{\&} to mark where lines align, and \texttt{\textbackslash{}\textbackslash{}} to end a line.

\begin{lstlisting}[language=latex]
\begin{align}
  r &= \frac{s}{c + v} \\
    &= \frac{e v}{c + v} \\
    &= \frac{e}{1 + c/v}.
\end{align}
\end{lstlisting}

Each line gets a number, and \texttt{align*} removes all of them. \texttt{\textbackslash{}nonumber} removes one.

\section{Maths as your own commands}\label{maths-as-your-own-commands}

If a symbol means something in your project, name it once, and every use will follow.

\begin{lstlisting}[language=latex]
\newcommand{\rate}{\ensuremath{r}}
\end{lstlisting}

\section{Tips}\label{tips}

\begin{itemize}
\tightlist
\item
  Spell out what each symbol means in the text around the equation.
\item
  Treat an equation as part of a sentence, with punctuation.
\item
  Don\textquotesingle t use \texttt{\$} or \texttt{\$\$} for displayed equations in LaTeX. Use \texttt{\textbackslash{}{[}\ \textbackslash{}{]}}, because \texttt{\$\$} skips LaTeX\textquotesingle s spacing rules.
\end{itemize}

